\part{Reasonless Intermediate Tokens}
\label{part:tokens}

\chapter{Formal Reconstruction of the Source Experiments}
\label{ch:valmeekam}

\section{The Domain}

The source experiments \cite{valmeekam2026} concern maze solving. A task $x$ is a maze with a start and a goal, serialized as a token sequence. The reference algorithm is A*, which maintains a priority queue of frontier cells ordered by the sum of path cost so far and a heuristic estimate of remaining cost, and which repeatedly expands the cell of least priority. Execution of A* on a maze induces a deterministic sequence of operations, namely the expansion of a cell and the insertion of its neighbors with their computed costs. A trace serialization writes this sequence as text. A formal validator $\Valid_{\mathcal S}$ parses a candidate trace and checks, operation by operation, that it is the sequence the algorithm would perform on the given maze, and that the final answer is the path recovered from the final state. The specification $\mathcal S$ is thus the execution semantics of A* together with the serialization grammar.

\section{Three Training Regimes}

Three training regimes are compared. In the solution-only regime the model is trained to emit the final path directly. In the correct-trace regime it is trained to emit the A* trace of the given maze followed by the path. In the swapped-trace regime each training example pairs a maze with the A* trace of a different maze and the correct path for the original maze, so that every trace is a syntactically well formed A* trace that is nonetheless not the trace of the maze it accompanies. The swapped regime is the experimental device that separates the form of a trace from its semantic relation to the task: it preserves the grammar, the vocabulary, the typical length, and the marginal statistics of traces while destroying their correspondence with the problem.

Evaluation is carried out on several maze distributions. Cross-distribution evaluation is informative because a model that has learned the procedure would be expected to transfer across distributions in which the procedure is the same, while a model that has learned distribution-specific regularities would be expected to degrade. In the source study the training mazes are generated by Wilson's algorithm in one setting and by a cave-style generator (SF-style) in another, and evaluation covers several generators including Kruskal, depth-first search, the Drunkard's Walk, and SF-style mazes, with differences in transfer reported between the regimes.

\section{Reported Findings}

The findings relevant to this monograph, stated as findings of the source study, are three. First, models trained on swapped traces reach high final-answer accuracy, and on some evaluation distributions exceed models trained on correct traces, in particular for the SF-style models on Drunkard and SF-style mazes. Second, trace validity as measured by the formal validator and answer accuracy diverge: models can produce correct answers with invalid traces. Third, in some post-training experiments with group relative policy optimization, answer accuracy improved while trace validity deteriorated. A further finding concerning partially swapped datasets is discussed in Chapter~\ref{ch:paradox}.

\section{The Precise Conclusion and Its Limits}

The conclusion supported is the following.

\begin{quote}
In the evaluated settings, high final-answer accuracy does not require the model to generate a trace that is valid under the prescribed A* execution semantics.
\end{quote}

The conclusion is a negative existence claim about a necessary condition, and its strength should not be inflated. It does not follow that the models performed no meaningful internal computation, because the answers were correct and something computed them. It does not follow that intermediate tokens never carry procedural content in other models or tasks. It does not follow that the traces had no causal influence on the answers, which is an interventional question the training-regime comparison does not address. What follows is that the semantic validity of the trace is not what accounts for the accuracy, and the explanatory burden shifts to whatever else does.

\section{What the Validator Checks}

The formal validator is the instrument on which the central finding rests, and its construction determines what the finding means. A validator for A* traces must establish several things for each operation recorded in the trace: that the cell expanded is the one with least priority in the queue as it stands after the preceding operations, that the neighbors inserted are the neighbors of that cell that are not walls, that the costs assigned to them are the correct sums of path cost and heuristic, and that ties are broken in the manner the specification prescribes. The last requirement is a source of subtlety, because A* as an algorithm does not fix a tie-breaking rule, and a validator must choose one. A trace that is a correct execution of A* under one tie-breaking rule may be rejected by a validator that assumes another. The reported validity rates are therefore rates of conformity to a particular serialization and tie-breaking convention, and an analysis that reads them as rates of algorithmic correctness in a broader sense should state the convention.

A second subtlety concerns what the validator does with a trace that is correct up to a point and then deviates. A binary validator rejects the whole trace. A finer validator can report the position of the first deviation, the fraction of operations that are correct, or the longest correct prefix. These measures are not interchangeable, and a model that produces mostly correct traces with occasional errors differs from one that produces traces that are wrong from the outset. The binary measure is the most conservative and the least informative, and the graded measures are preferred where the object of study is how the traces degrade.

\section{The Swapped Regime as an Intervention on Supervision}

The swapped regime has the logical structure of a controlled intervention on the training data. It holds fixed the marginal distribution of traces, because the set of traces used is the same as in the correct regime, and holds fixed the marginal distribution of mazes and of answers. It breaks the dependence between the maze and its trace, replacing it with independence. The dependence between the maze and the answer is untouched. The comparison between the correct and swapped regimes therefore isolates the contribution of the trace-maze dependence to the performance of the trained model, with everything else held constant, and a performance difference between the regimes is attributable to that dependence under the usual assumptions of a controlled experiment.

The contrast with the solution-only regime has a different logical structure. The solution-only model has no trace at all, so its comparison with the swapped model isolates the effect of having additional token positions with a particular marginal distribution of content. If the swapped model outperforms the solution-only model, the additional positions are useful even when their content is independent of the task during training. This is the comparison most relevant to the scaffolding hypothesis. If the swapped model does no better than the solution-only model, the utility of the correct trace would have to lie in its dependence on the maze, which the swapped regime has removed.

\section{Reading Cross-Distribution Results}

Evaluation on distributions other than the training distribution tests whether what the model learned is specific to the training mazes. A model that has learned the procedure should transfer, to the extent that the procedure is the same on the new mazes and the model can represent the new inputs. A model that has learned heuristics adapted to the training distribution, such as typical corridor shapes, should degrade. The inference from transfer to procedure is weaker than it appears, however, because transfer can also occur when the new distribution shares statistical structure with the old one, and the degree of shared structure is not controlled by the experimenter unless the distributions are constructed to vary it. The mazes described in the source as Drunkard-style and SF-style differ in their generation processes, and the extent to which they differ in the features relevant to a learned heuristic is an empirical matter that the performance numbers alone do not settle.

\section{Threats to Inference}

Four threats limit what can be concluded from the experimental design. The first is scale: the models are small and trained on a narrow task, and results for such models need not hold for large pretrained models whose representations are shaped by far broader data. The second is the choice of task: maze solving by A* is a deterministic algorithmic problem with a clean specification, and the dissociation observed there may not carry over to tasks with weaker specifications, such as open-ended mathematics or planning in an uncertain world. The third is the choice of specification: validity is conformity to one serialization of one algorithm, and other characterizations of what the model does might classify the same traces differently. The fourth is the training procedure: the results were obtained under particular optimizers, learning rates, and data volumes, and the dependence on these is not known from a single study. None of the threats undermines the narrow conclusion stated above, which is a statement about the evaluated settings. They bear on how far the conclusion may be extended, and the monograph extends it only as a hypothesis for further testing.

\chapter{Information, Computation, and Intermediate Tokens}
\label{ch:info}

\section{Problem Information and Computational Transformation}

A distinction between information about the problem and transformation of that information is needed to state what intermediate tokens can and cannot contribute. Let $x$ denote the task, $z$ the intermediate sequence, and $y$ the final answer. An autoregressive model factors the joint distribution as
\[
p_\theta(z,y\mid x)=p_\theta(z\mid x)\,p_\theta(y\mid x,z).
\]
Let $y^\ast=f(x)$ denote the answer defined by the task. Because $z$ is generated from $x$ by the model, possibly with independent sampling noise, the variables form the Markov chain $y^\ast\to x\to z$, and the data processing inequality gives $I(y^\ast;Z)\le I(y^\ast;X)$. The sequence $z$ therefore cannot contain more information about the correct answer than $x$ already contains. The inequality does not say that $z$ cannot help. It says that any benefit from $z$ is a benefit of computation and not of new external information: $z$ can make information already present in $x$ available in a form that the next-token computation can use within its fixed depth.

\section{Training Traces and Generated Traces}

A distinction between two sources of $z$ matters for interpretation. In training, a trace is a supervised target. In the swapped regime it is statistically independent of the specific maze it is paired with, given the maze distribution. At inference, the trace is generated by the model from the input and is therefore statistically dependent on the input even when it is procedurally invalid. A generated invalid trace can encode functions of the input. This is the reason that the swapped-regime finding does not imply the traces are uninformative at inference: the training supervision was uninformative, but the learned generator is not required to reproduce that uninformativeness.

\section{Extra Tokens as Computational Opportunity}

A transformer of fixed depth performs a bounded amount of serial computation per token. Additional tokens provide additional passes, with the intermediate results held in the sequence that later positions attend to. In this respect intermediate tokens resemble scratchpads \cite{nye2021}, iterative algorithms that reuse a state, recurrent computation, and learned latent workspaces. The resemblance is structural and should not be taken to imply equivalence. A scratchpad whose contents are semantically meaningful differs from one whose contents are an uninterpreted code, even if both give the network additional serial steps. The comparison that matters is between sequences matched on length and token budget and differing in semantic relation to the task, which is the comparison the swapped regime provides.

\section{The Serial Depth Argument}

A transformer with $L$ layers performs at most $L$ sequential applications of its layer stack to the information at any given position, and the information available to a position is limited to what earlier positions have computed and exposed through attention. When the model produces a token and that token is fed back as input, the next position can use the result of the preceding computation as though it were part of the input, and the effective serial depth of the whole computation increases by up to $L$. For a sequence of $k$ generated tokens the serial depth available is of the order of $kL$, instead of $L$. Problems whose solution requires serial composition of more than $L$ dependent operations cannot be solved in a single forward pass, under standard assumptions about what a bounded-depth network computes, and can be solved with a sufficient number of intermediate tokens that carry the intermediate results.

The argument establishes that additional tokens are useful for a reason that has nothing to do with their semantic content, namely that they provide additional serial steps. It also shows what is required of them: to carry intermediate results forward, the tokens must encode those results in a way that later positions can decode. This is a requirement on the code, which need not be the code of the prescribed algorithm. A model that has learned its own code for intermediate results can use tokens as a scratchpad whose entries are meaningful to itself and not to a reader who applies the A* specification. The swapped-trace result is consistent with this reading, in which the learned code diverges from the prescribed one, and it is also consistent with readings in which the tokens carry no intermediate results at all. The serial depth argument by itself does not distinguish these.

\section{Exposure and Self-Conditioning}

During training with supervised traces, the model is conditioned on the correct previous tokens at every position, a procedure known as teacher forcing. At inference, the model is conditioned on its own previous outputs. If its outputs contain errors, later positions see inputs unlike any they were trained on, and the effect of an error can compound. This mismatch between training and inference conditions has two relevant consequences. First, a model trained on correct traces may produce traces whose later portions are degraded by earlier errors, and such traces will fail a strict validator while the answer, if it depends on a computation less sensitive to the errors, remains correct. Second, a model trained on swapped traces never sees traces that correspond to its inputs, and so cannot have learned to rely on them for the content of the answer. It must learn to compute the answer from the input by some route that does not require the trace to be a reliable record. Whether it then uses the trace at all is an empirical question that the interventional comparison of Chapter~\ref{ch:scaffolding} addresses.

\section{Entropy of Generated Traces}

The statistical dependence of generated traces on the input can be quantified. For a trained model, the conditional entropy $H(Z\mid X)$ of the generated trace given the task measures how much of the trace is determined by the task, and the mutual information $I(Z;X)=H(Z)-H(Z\mid X)$ measures how much information about the task the trace carries. A model that has learned a procedure produces a trace with high mutual information with the task, because the trace is a deterministic function of it. A model that emits a task-independent filler produces a trace with mutual information near zero. Generated traces of models trained on swapped data are an empirical object whose mutual information with the task is not known in advance, and measuring it distinguishes the case in which the model has learned to produce task-dependent invalid traces from the case in which it produces task-independent text. The first would suggest a learned code, and the second a pure serial-computation affordance.
This measurement can be carried out with standard estimators on the generated traces and does not require access to the model's internal states.

\chapter{The Computational Scaffolding Hypothesis}
\label{ch:scaffolding}

\section{Statement}

\begin{hypothesis}[Computational scaffolding]\label{hyp:scaffold}
Intermediate tokens can function as computational scaffolds whose utility is not reducible to their apparent procedural semantics. The utility depends on properties of the sequence, such as its length, its statistical regularity, and the conditional structure the model has learned to associate with it, that are partly or wholly independent of whether the sequence is a valid derivation.
\end{hypothesis}

The hypothesis is deliberately weaker than the geometric hypothesis of Part~\ref{part:faithfulness}. It asserts only that semantic validity is not the sole source of utility, and it is compatible with several mechanisms, including provision of additional serial computation, stabilization of an internal state through repeated conditioning, and priming of a mode of processing.

\section{An Interventional Measure of Utility}

Utility is defined by intervention. Let $z_0$ be a specified control sequence and let $\doo(Z=z)$ denote controlled replacement of the generated prefix by $z$. The utility of $z$ on task $x$ relative to $z_0$ is
\[
\Delta U(z;x)=\Pr\big(\Corr=1\mid x,\doo(Z=z)\big)-\Pr\big(\Corr=1\mid x,\doo(Z=z_0)\big).
\]
This quantity depends on the intervention protocol and on the control, and it is not an intrinsic measure of the semantic value of $z$. Two different controls, such as an empty prefix and a prefix of random tokens of equal length, answer different questions. The first asks whether any prefix helps. The second asks whether the specific content of the prefix helps beyond occupying positions.

\section{What Scaffolds Might Depend On}

The hypothesis suggests a structured investigation of what a useful scaffold depends on. Candidates are length, syntax, domain-specific vocabulary, statistical regularity of the token sequence, and learned conditional structure, meaning the dependence of later tokens on earlier ones as the model has learned it. Each can be varied while the others are held fixed. Repeated tokens preserve length and vocabulary while removing structure. Random tokens preserve length while removing vocabulary and structure. Shuffled valid traces preserve vocabulary and marginal structure while removing the dependence on the task. Learned prefixes, optimized by gradient descent against task success with no constraint on form, preserve nothing but the optimization target and indicate how much utility is available at all.

\section{Proposed Comparisons}

The comparison suite is a factorial design of prefix type against matched computational budget. The prefix types are correct traces, swapped traces, random tokens, repeated tokens, and learned prefixes, each evaluated at token counts matched to those of the correct traces. The outcome is answer accuracy on held-out mazes from each evaluation distribution, and the analysis estimates $\Delta U$ for each type relative to the empty prefix and relative to the random prefix. A finding that swapped traces and random tokens give equal utility at equal length would favor a pure serial-computation account. A finding that swapped traces outperform random tokens would indicate that the learned conditional structure contributes. The design is specified in full in Appendix~\ref{app:protocols}.

\section{A Taxonomy of Mechanisms}

Several mechanisms can account for a case in which a trace is useful and invalid, and they differ in what they predict. A serial-computation mechanism holds that the tokens supply extra forward passes, and that their content is incidental provided the length is sufficient. It predicts that random tokens of the same length are as useful as swapped traces and that repeated tokens are nearly so. A learned-code mechanism holds that the model has developed its own encoding of intermediate results in the token stream, so that the tokens carry information used by later positions. It predicts that swapped traces are more useful than random tokens, that the tokens have high mutual information with the task, and that replacing a generated trace with another trace from the same model on a different maze lowers accuracy. A priming mechanism holds that the tokens place the model in a mode of processing, much as a task description does, and that their specific content matters little beyond identifying the mode. It predicts that a short prefix of the right type suffices and that length beyond some threshold adds nothing. A distribution-matching mechanism holds that the trace keeps the model's inputs within the range of its training distribution, so that the later computation operates in a familiar regime. It predicts that traces from the model's own distribution are useful and that out-of-distribution traces, even if valid, are harmful.

These predictions are distinguishable by the factorial design of the preceding section, supplemented by measurements of mutual information and by cross-maze replacement. No single mechanism is expected to account for everything, and the mixture of mechanisms may differ across models and tasks. The relevance for the geometric account is that each of them is compatible with a description in terms of trajectories: serial computation lengthens the path, a learned code shapes it, priming sets its starting region, and distribution matching keeps it near familiar states. The geometric account is a way of describing the effects and does not by itself choose among the mechanisms.

\section{Relation to Soft Prompts and Learned Prefixes}

The learned-prefix arm of the design has a close relative in the work on prompt tuning, in which continuous vectors prepended to the input are optimized by gradient descent to steer a frozen model toward a task \cite{lester2021}. Such prefixes have no semantic interpretation as text, and yet they can be as effective as full fine-tuning in some regimes. This is a demonstration that the utility of a prefix need not depend on its being interpretable as language. The difference from the intermediate tokens studied here is that soft prompts are continuous and are optimized with the model frozen, whereas intermediate tokens are discrete and the model is trained to produce them. The existence of effective uninterpretable prefixes supports the plausibility of the scaffolding hypothesis and does not establish it for generated discrete tokens.

\section{Statistical Considerations}

The comparisons proposed involve differences in accuracy that may be small, evaluated on models trained with different random seeds, and the analysis must respect that structure. Variation across seeds can be as large as the differences between conditions, so conclusions require several seeds per condition and an analysis that treats the seed as the unit of replication for claims about training, and the maze as the unit for claims about evaluation. The comparisons among five prefix types and several evaluation distributions involve many hypothesis tests, and a correction for multiple comparisons or a hierarchical model should be used. Effect sizes should be reported with intervals, and a failure to detect a difference should be reported with the interval that shows how large a difference the experiment could have detected, since a null result from an underpowered comparison does not support equivalence.

\chapter{The Generalization Paradox}
\label{ch:paradox}

\section{The Observation}

A surprising feature of the source results is that training on correct algorithmic traces can generalize less well to some distributions than training on mismatched ones. The observation is surprising only against the background expectation that supervision closer to a correct algorithm should yield behavior closer to it. The source study also reports a U-shaped performance trend across partially swapped datasets, in which accuracy is not monotone in the fraction of swapped traces. The shape is reported as an observation, and this chapter treats its explanation as open.

\section{Competing Explanations}

Four families of explanation can be distinguished, and they make different predictions. Overconstraint holds that correct traces impose a specific serialization of the search that the model must reproduce exactly, and that errors in reproduction propagate, whereas a model free of such constraint can compute the answer by whatever route it finds. Imitation of incidental details holds that correct traces contain regularities particular to the training distribution, such as typical expansion orders in a given maze style, that the model learns and that fail to transfer. Conflicting supervision holds that when the trace and the answer are produced by the same network under a single loss, the demands of faithfully reproducing the trace and of computing the answer pull in different directions. Optimization dynamics holds that the differences are properties of the loss landscape reached from the different data and not of any representational difference. The families are not exclusive.

\section{Competence Versus Task Completion}

Genuine algorithmic generalization should be distinguished from successful task completion on unfamiliar input distributions. A model that solves unfamiliar mazes has generalized in the weak sense that its answers are correct. It has generalized in the strong sense only if what it does on the new distribution is the same computation that it does on the old. The two senses coincide for a model that implements the algorithm and come apart for a model that implements a collection of distribution-specific heuristics that happen to succeed. Evidence for the strong sense requires examining the internal computation, and the geometric measurements of Chapter~\ref{ch:hidden-states} are one candidate.

\section{An Experimental Design}

To separate algorithm-specific imitation from broader computational competence, train models in each regime on a family of maze distributions that differ along controlled axes, evaluate on held-out distributions that vary along each axis separately, and measure both answer accuracy and the agreement of internal trajectories across distributions, for instance by the similarity of the sequences of basins traversed. If the correct-trace regime shows lower answer transfer but higher internal agreement, algorithm-specific imitation without competence on novel structure is indicated. If the swapped regime shows higher answer transfer with lower internal agreement, success by distribution-specific routes is indicated.

\section{A Stylized Model of the Non-Monotone Pattern}

A simple illustrative model shows how a non-monotone dependence on the fraction of swapped traces can arise without any exotic mechanism. The model is offered only to show that the pattern is not paradoxical, and it makes no claim about the source experiments. Suppose a model can reach an answer by two routes. The first route uses the trace as a record of intermediate results and is learned to the extent that the training traces are correct, so its strength grows with the fraction of correct traces. The second route computes the answer without relying on the trace and is learned to the extent that the trace is unreliable, so its strength grows with the fraction of swapped traces. At a fraction near zero swapped, the first route is strong and the second weak, and the model performs well on the training distribution through the first route. At a fraction near one, the second route is strong and the first weak, and the model performs well through the second. At intermediate fractions, neither route is fully developed, because the supervision for each is contaminated by the other, and the model may be left with two weak routes and a conflicted treatment of the trace. Performance may then dip, producing a U-shaped dependence.

The model makes testable predictions. It predicts that at intermediate fractions the answer depends on the trace to an intermediate degree, so that the interventional influence of the trace varies monotonically with the fraction even though accuracy does not. It predicts that the internal trajectories at intermediate fractions show a mixture of the signatures of the two routes. These predictions can be checked against the interventional and geometric measurements proposed in the monograph, and they distinguish the stylized model from alternatives in which the dip arises from optimization difficulty without any change in routes.

\section{What Each Explanation Predicts}

The four families of explanation separate under suitable measurement. Overconstraint predicts that models trained on correct traces fail primarily by compounding errors in the trace, so that failures cluster in long traces and in the later portions of traces, and that giving the model the correct prefix at inference restores its accuracy. Imitation of incidental details predicts that models trained on correct traces show distribution-specific regularities in their traces, such as characteristic expansion orders, which transfer poorly, and that accuracy on held-out distributions correlates with the degree to which the trace departs from the training distribution's regularities. Conflicting supervision predicts that the loss on the trace and the loss on the answer trade off during training, which can be seen in the training curves, and that weighting the answer loss more heavily improves transfer. Optimization dynamics predicts that the differences between regimes diminish when the optimizer and schedule are tuned separately for each regime, and persist only under a fixed shared configuration. These are cheap tests relative to the experiments of the program, and they should be run first, since a positive result for any of them changes the interpretation of the swapped-trace advantage.

\chapter{Trace Length and the Appearance of Adaptive Deliberation}
\label{ch:length}

\section{The Comparison}

The source study compares the lengths of generated traces with the lengths of reference A* traces on the same mazes. A positive relation between generated length and a measure of problem difficulty is often read as evidence that the model allocates more deliberation to harder problems. The reading requires care at each step.

\section{Measures of Association}

The Pearson correlation measures linear association and is sensitive to outliers. Rank correlation measures monotone association and is robust to nonlinear but monotone relations. Neither detects non-monotone relations, and both are distorted by censoring when generated traces are truncated at the context limit, because the longest traces are then recorded at the limit regardless of what they would have been. An analysis of length should report the association by more than one measure, state the censoring, and examine the relation graphically.

\section{Difficulty Is Relative to an Algorithm}

The number of expansions that A* performs on a maze is an algorithm-relative measure of difficulty. A different algorithm, or the same algorithm with a different heuristic, yields a different count on the same maze. It is not an intrinsic measure of computational complexity and should not be treated as one. A model whose trace length tracks A* expansion count is tracking a property of the reference algorithm, which could reflect either the use of something like that algorithm or the correlation of both quantities with some third property of the maze, such as its size or branching.

\section{Four Distinct Quantities}

Four quantities should be defined separately. Task difficulty is a property of the task relative to a stated solver. Computational expenditure is the amount of computation actually performed by the model, which is not directly observed. Generated sequence length is the observed count of tokens. Problem-adaptive resource allocation is a relation between difficulty and expenditure, namely that expenditure increases with difficulty beyond what is explained by input size. Sequence length is an observable proxy for expenditure, and expenditure is an unobserved quantity that the proxy approximates only to the extent that tokens correspond to units of computation.

\section{Tests for Adaptive Allocation}

Tests for whether a model allocates computation according to uncertainty, task structure, or distributional familiarity can be built from controlled variation. Hold input size fixed and vary solution structure to test the dependence on task structure. Present tasks from a familiar distribution and matched tasks from an unfamiliar one to test familiarity. Measure the model's own uncertainty over the answer at the point where the trace begins, and test whether trace length is predicted by that uncertainty after controlling for size and structure. Each test isolates one of the three candidate bases of allocation and a model may show more than one.

\section{Censoring Formalized}

Let $L^\ast$ be the length the model would generate without a limit and let $c$ be the context limit. The observed length is $L=\min(L^\ast,c)$. For tasks whose natural length exceeds $c$, the observed length is $c$ regardless of how much longer the trace would have been, and the observed data have a point mass at $c$. A correlation between $L$ and a difficulty measure computed on such data underestimates the association that would hold for $L^\ast$, because the censored observations compress the upper range. Rank correlations suffer from the ties at $c$, and the standard remedy for ties changes their values. Methods from survival analysis, which treat $c$ as a censoring time and estimate the distribution of $L^\ast$ accordingly, are appropriate when the proportion of censored observations is substantial, and the analysis should in any case report that proportion.

\section{Confounding by Input Size}

A relation between trace length and a difficulty measure is easily produced by a common dependence on the size of the input. Suppose both quantities increase with the size $s$ of the maze, so that $L=as+\epsilon_1$ and $D=bs+\epsilon_2$ with independent noise, and with no direct relation between the trace length and the difficulty at a given size. The correlation between $L$ and $D$ is then positive and may be strong, since both track $s$, although the model allocates no computation in response to difficulty beyond what size already implies. The partial correlation of $L$ and $D$ given $s$ removes the effect of the common cause and is zero in this construction. A claim of adaptive allocation therefore requires the partial association, controlling for size and for other task features that determine length mechanically. The reported association between generated length and A* expansion count should be read in this light, and the proposals of this chapter specify the controls that would convert it into evidence of adaptivity or show that it is absent.

\section{Role in the Argument}

The analysis of length bears on the main argument in two ways. The first concerns the scaffolding hypothesis. If generated length tracked problem difficulty after the controls above, the tokens would be functioning as a resource allocated in response to need, which would favor accounts in which the content of the trace is doing work. The source study reports a negligible correlation between generated and reference trace lengths for models trained on Wilson mazes, and a pattern for other distributions that depends on the measure and the maze family \cite{valmeekam2026}. A length that does not follow difficulty is what the serial-computation and priming mechanisms of Chapter~\ref{ch:scaffolding} would predict, since both treat the tokens as a roughly fixed affordance and not as an adaptive resource. The length analysis therefore supplies evidence that discriminates among the mechanisms, and its place in the argument is as a test of one of their predictions.

The second concerns the structure of the promotions of Chapter~\ref{ch:authority}. The reading of trace length as deliberation is itself a promotion: an observable property, the number of tokens, is taken to establish an unobserved one, the amount of effort spent on the problem. A longer trace is read as more thought, and the reading flatters the system and informs the observer's trust. Separating length from expenditure and expenditure from adaptive allocation is the same discipline applied to a further case, and it shows that the appearance of deliberation can be as misleading as the appearance of derivation.
